% ----------------------------------------------------------
% Introdução
% ----------------------------------------------------------
\chapter{Introdução}

A gestão acadêmica de instituições de ensino superior, envolve processos complexos que demandam não apenas questões pedagógicas, mas também precisão administrativa e documental. No contexto do Instituto Federal de Pernambuco (IFPE), um dos documentos mais importantes (talvez o mais importante) para a criação e manuntenção de cursos é o Projeto Pedagógico de Curso (PPC). Este documento define toda a estrutura acadêmica, desde a justificativa, objetivos, matriz curricular com suas respectivas cargas horárias, ementas e bibliografias.

\section{Problema}
Atualmente, a elaboração de um PPC é um trabalho bastante dispendioso, realizada na maior parte do tempo de forma manual através de editores de texto e planilhas eletrônicas desvinculadas, além de ser modificado por diversas pessoas ao longo do tempo. Esse processo manual expõe a gestão acadêmica a uma série de riscos: erros de cálculo no somatório de cargas horárias, inconsistência entre matriz curricular, ementas e anexos, formatação fora dos padrões institucionais e, sobretudo, um consumo excessivo de tempo de coordenadores e membros dos Núcleo Docentes Estruturantes (NDE). A complexidade aumenta muito ao considerar as atualizações nas diretrizes curriculares nacionais, na lista de cursos tecnólogos catalogados no MEC e a necessidade de padronização entre os diversos \textit{campi}.
A falta de padronização dos PPCs de diferentes cursos dos Institutos Federais é uma característica notável da ausência de ferramentas específicas para esta finalidade reproduzida diversas vezes por diferentes IFs, e resulta em retrabalho e na dificuldade de manter a integridade dos dados pedagógicos. Diante deste cenário, torna-se necessária a adoção de soluções tecnológicas capazes de gerenciar a entrada de dados e a geração dos documentos, garantindo a conformidade com as normas vigentes e permitindo que a gestão foque na qualidade pedagógiva, e não na formatação documental.

\section{Solução}
Neste contexto, este trabalho apresenta o desenvolvimento do \textbf{xPPC}, uma aplicação \textit{desktop} construída sobre a plataforma Java, projetada para automatizar a criação e padronização dos Projetos Pedagógicos de Curso no âmbito do IFPE. Utilizando tecnologias mordernas de engenharia de software e bibliotecas de manipulação de documentos, o sistema propões uma abordagem onde o usuário insere os dados estruturados em formulários intuitivos e aplicação se responsabiliza por realizar cálculos, validar informações e gerar o documento final devidamente formatado, mitigando erros humanos e agilizando o fluxo admnistrativo da instituição. 

\section{Objetivos}
\subsection{Objetivo Geral}
\subsection{Objetivos específicos}

\section{Justificativa}
\section{Organização do trabalho}